\documentclass[11pt]{article}
\usepackage[margin=1in]{geometry}
\usepackage{hyperref}
\usepackage{listings}
\usepackage{xcolor}
\usepackage{amsmath}

\lstset{
  basicstyle=\ttfamily\small,
  breaklines=true,
  columns=fullflexible,
  frame=single,
  backgroundcolor=\color{gray!10},
  showstringspaces=false
}

\title{URI OPeNDAP Test Case: Working Log (Prompt \#2)}
\author{Compiled for P. Cornillon}
\date{2026-05-16}

\begin{document}
\maketitle

\section*{Purpose}

Working log for Prompt~\#2 of the URI OPeNDAP test case. Findings
themselves are in \texttt{docs/uri\_test\_case\_2.tex}; this file
records process, command lines, decisions, and corrections to
Prompt~\#1.

\section{Prompt \#2 (2026-05-16 19:42 UTC)}

\subsection{Instruction}

\emph{``Generate a table for Readable Branch \#1: SST\_Orbits/. For
each variable in the dataset, indicate whether the semantic
(COARDS) metadata essential for using that variable are complete.
If the metadata are complete, state that explicitly. If they are
incomplete, list the missing COARDS metadata and provide your best
guess for what the values should be. In addition, identify, again
for each variable, any semantic metadata that are not strictly
required by COARDS but that you believe would significantly
facilitate use and interpretation of the dataset.''}

\subsection{Re-reading Prompt \#1 outputs}

The two files produced for Prompt~\#1
(\texttt{docs/uri\_test\_case\_1.tex},
\texttt{docs/test\_case\_uri\_log\_1.tex}) describe a Hyrax server
with seventeen top-level containers, of which only
\texttt{SST\_Orbits/} and \texttt{gradients\_by\_period/} are
publicly readable. The orbit-file branch was characterized as
``richly CF-annotated'' but two specific claims in those documents
were wrong and are corrected below.

\subsection{Corrections to Prompt \#1}

\begin{enumerate}
  \item Prompt~\#1 said the orbit DMR exposes no root-group
    (global) attributes. \textbf{False.} A more careful parse of
    the DMR (using a depth-aware regex walker, not a flat one)
    finds 22 root-group attributes including
    \texttt{title="Conditioned 1km (L2) MODIS AQUA SST"},
    \texttt{Conventions="CF-1.5"}, \texttt{institution="University
    of Rhode Island - Graduate School of Oceanography"},
    \texttt{creator\_name="Peter Cornillon"},
    \texttt{project="SST Fronts"}, and a full ACDD-style discovery
    block. The earlier omission was an artifact of my flat-regex
    DMR parser, which was treating attribute elements as if they
    only appeared inside variable elements.
  \item Prompt~\#1 said the Int32 \texttt{latitude} and
    \texttt{longitude} variables in the orbit file declare
    \texttt{units=degrees\_*} but do \emph{not} declare a
    \texttt{scale\_factor}. \textbf{False.} Both variables do
    carry \texttt{scale\_factor=0.001}, \texttt{add\_offset=0},
    and corresponding \texttt{valid\_min/max} in the DMR. The
    Int32 raw values \emph{do} need to be rescaled by $10^{-3}$
    --- that conclusion was right --- but the file does say so
    explicitly. Prompt~\#1's analysis script was failing to
    capture more than the first few attributes per variable.
\end{enumerate}

I record these so that anything downstream that quotes Prompt~\#1
on those two points knows to discard the claim.

\subsection{Approach for Prompt \#2}

\begin{enumerate}
  \item Re-fetch the DMR for the representative orbit file
    \begin{lstlisting}
https://sst-aqua.gso.uri.edu/opendap/SST_Orbits/2024/01/
  AQUA_MODIS_orbit_115220_20240101T011558_L2_SST-URI_24-2.nc4.dmr.xml
    \end{lstlisting}
    with a depth-aware parser
    (\texttt{scripts/sst\_orbits\_dmr\_inventory.py}) that walks
    every \texttt{<Group>} and emits, for every variable, every
    attribute as \texttt{(name, type, value)}.
  \item Score each variable against a deliberately strict reading
    of COARDS:
    \begin{itemize}
      \item \texttt{units} present and UDUNITS-compliant;
      \item \texttt{long\_name} present;
      \item \texttt{\_FillValue} present (for non-scalar,
        non-string data);
      \item if either of \texttt{scale\_factor},
        \texttt{add\_offset} is set, the other must be too;
      \item if the variable is a latitude / longitude / time
        axis, the \texttt{units} string must match the
        COARDS-recognized form (e.g.\
        \texttt{degrees\_north}; for time, \texttt{<unit> since
        <reference>}).
    \end{itemize}
    The scorer is in
    \texttt{scripts/sst\_orbits\_coards\_audit.py}; results are
    persisted to \texttt{/tmp/sst\_audit.json} so that the
    findings table can be regenerated without re-fetching.
  \item For each variable, also collect \emph{non-COARDS}
    recommendations (CF/ACDD-style attributes that would help
    interpretation: \texttt{coordinates},
    \texttt{ancillary\_variables}, \texttt{cell\_methods},
    \texttt{flag\_values/flag\_meanings}, \texttt{axis},
    \texttt{calendar}, \texttt{source}, instrument provenance
    notes).
  \item Compose the assessment into
    \texttt{docs/uri\_test\_case\_2.tex} as three per-group
    tables (root, \texttt{/Regrid\_to\_L2eqa},
    \texttt{/contributing\_granules}).
\end{enumerate}

\subsection{Commands}

\begin{lstlisting}
# 1. Pull the DMR XML and parse it
python3 scripts/sst_orbits_dmr_inventory.py > /tmp/sst_inv.json

# 2. Run the COARDS scorer + recommendation generator
python3 scripts/sst_orbits_coards_audit.py
# -> writes /tmp/sst_audit.json
\end{lstlisting}

\subsection{Headline numbers}

Of the 44 atomic variables across the three groups of the orbit
file, the strict scorer marks \textbf{21 Complete} and
\textbf{23 Incomplete} for COARDS. The breakdown by group:

\begin{lstlisting}
/                              11 / 20 complete
/Regrid_to_L2eqa                9 / 17 complete
/contributing_granules          1 /  7 complete
\end{lstlisting}

The largest single contributor to ``Incomplete'' status is the use
of the units string \texttt{"C"} (and \texttt{"C/km"}) where
UDUNITS-compatible code would expect \texttt{degree\_Celsius}
(\texttt{degree\_Celsius km-1}). That single substitution would
close out 15 of the 23 incomplete entries. The remaining 8 are:
two \texttt{region\_start/end} arrays with no units at all, three
\texttt{contributing\_granules} index arrays with no units, two
\texttt{start\_time/end\_time} arrays whose \texttt{units="seconds"}
lack a reference epoch, and the scalar \texttt{DateTime} which has
the same problem.

\subsection{Limitations of the audit}

\begin{itemize}
  \item The scorer treats \texttt{long\_name} as required even
    though COARDS only \emph{recommends} it. This makes the
    threshold a little stricter than the literal standard but
    keeps the table useful as a checklist.
  \item Some files in the SST\_Orbits/ branch may have been
    re-processed since this sample (the sample's
    \texttt{date\_created} is 2025-05-17). Different epochs may
    carry slightly different attribute sets; everything below
    refers to the specific sample file unless otherwise noted.
  \item The audit does not look at \emph{values} except for unit
    mislabellings I noticed (\texttt{L2eqa\_AMSR\_E\_wind\_speed}
    declared \texttt{units="mm"} while its packing is consistent
    with $\text{m s}^{-1}$;
    \texttt{L2eqa\_MODIS\_num\_SST} declared \texttt{units="C"}
    while it is plainly a count). A proper audit would compare
    sampled values against the declared range.
\end{itemize}

\subsection{Session timing}

Started 2026-05-16 19:42\,UTC. Closed log 2026-05-16
$\sim$20:05\,UTC.

\end{document}
